\documentclass[11pt]{article}
\usepackage{fontspec}
\usepackage[french,english]{babel}
\usepackage{amsmath}
\usepackage{amssymb}
\usepackage{amsthm}
\usepackage{mathrsfs}
\usepackage{tikz-cd}
\usepackage{enumitem}

\newcommand{\printedpage}[1]{\marginpar{\footnotesize p.~#1}}
\newcommand{\msp}[1]{\marginpar{\footnotesize [#1]}}
\newcommand{\Cat}{\mathrm{Cat}}
\newcommand{\Dl}{\boldsymbol{\Delta}}
\newcommand{\fl}{\longrightarrow}
\newcommand{\Ob}{\mathrm{Ob}\,}
\newcommand{\Hot}{\mathrm{Hot}}
\newcommand{\hot}{\mathrm{hot}}
\newcommand{\uHom}{\underline{\mathrm{Hom}}}
\newcommand{\id}{\mathrm{id}}
\newcommand{\pr}{\mathrm{pr}}

\begin{document}

\printedpage{--}
% head: (none; no running head and no page number printed on this page)

\begin{center}
{\Large \textsf{THE DERIVATORS}}

\bigskip

\textbf{ALEXANDRE GROTHENDIECK}

\bigskip\bigskip

\textbf{Chapter XII}

\medskip

\textbf{Characterization of $W_\infty$. $W$-proper,\\
$W$-smooth functors etc. Amalgamated sums and\\
$W$-cocartesian squares in $\Cat$.}
\end{center}

\clearpage

\printedpage{--}
% head: (none; no running head and no page number printed on this page)

This text was deciphered and transcribed in \LaTeXe{} by M. Künzer. It was edited by M. Künzer, J. Malgoire, and G. Maltsiniotis. The transcription is as faithful as possible to the manuscript. For the few obvious corrections, or rare comments by the editors, as well as for the original numbering of the manuscript pages, typewriter characters \texttt{[typewriter]} in brackets are used. A question mark in brackets means that one is not sure of the preceding word.

This edition is provisional. Remarks, comments and corrections are welcome. Send a message to:

\textbf{[e-mail address removed]}

\hfill G. Maltsiniotis\hspace*{2em}

\clearpage

\printedpage{1}
% head: A. Grothendieck : Dérivateurs, Ch. XII, Édition des Universités de Montpellier II et Paris VII
% foot: Version 7 décembre 2002

\msp{page 1}

\hfill 19.11.90

\section*{I\quad Characterization of $W_\infty$.}

I would like to characterize $W_\infty$ as a subset of $\mathrm{Fl}(\Cat)$, by its properties as a localizer, and without appealing to knowledge of notions already using $W_\infty$, such as properness etc. I would like at least some properties of a localizer $W$, familiar for $W_\infty$, which are so strong that they imply $W \subseteq W_\infty$. Of course, in the proofs I allow myself to use everything I know about $W_\infty$.

Of course, at least a form of saturation will be needed, perhaps a very weak one, I presume the following will suffice (\footnote{(no, I had to strengthen it in 2 [?]).}).

\begin{itemize}[leftmargin=*]
\item[\textbf{W(1)}] (Saturation axiom.) a) The identities are in $W$, b) if $f$, $g$ \texttt{[are]} composable and two among $f$, $g$, $gf$ are in $W$, the third too, c) if $fg \in W$, $gf \in W$, then $f, g \in W$.
\end{itemize}

A particularly crucial property, which already distinguishes $W_\infty$ from the set of homotopisms (a notion too strict in many respects), is the axiom I called L4 bis:

\begin{itemize}[leftmargin=*]
\item[\textbf{W(2)}] (Localization axiom.) Let $f : X \fl Y$ be an arrow in $\Cat/S$, such that $\forall\ s \in \Ob S$, $f_s : X/s \fl Y/s$ is in $W$ (we shall say that $f$ is in $W$ ``locally over $S$\,''). Then $f \in W$.
\end{itemize}

\texttt{[The axioms W(1) etc. are also called} \textcircled{1} \texttt{etc. in the original.]}

In the case $Y = S$, this gives

\textbf{Corollary:} \textit{If $f : X \fl Y$ is $W$-aspheric, i.e. $y\backslash X \fl y\backslash Y$} \texttt{[rather $X/y \fl Y/y$]} \textit{in $W$ for every $y \in \Ob Y$, then $f \in W$} (\footnote{One can consider the validity of the corollary as a weakened variant W(2 bis) of W(2) above. In proposition 1, one uses only W(2 bis), a weak form of W(2).}).

\textbf{Proposition 1.} \textit{If $W$ satisfies} W(1) \textit{and} W(2)\textit{, the following properties of $W$ are equivalent:}

\begin{itemize}[leftmargin=*]
\item[\textit{a)}] \textit{For every $X$ in $\Cat$, $X \times \Dl^1 \xrightarrow{\pr_1^X} X$ is in $W$ (i.e. $\Dl^1 \fl e$ is ``universally in $W$\,'').}

\item[\textit{b)}] \textit{If $X$ in $\Cat$ has a final object, then $X$ is $W$-aspheric, i.e. $X \fl e$ is in $W$.}

\item[\textit{c)}] \textit{Every homotopism is in $W$.}
\end{itemize}

\clearpage

\printedpage{2}
% head: A. Grothendieck : Dérivateurs, Ch. XII, Édition des Universités de Montpellier II et Paris VII
% foot: Version 7 décembre 2002

\textsc{Proof.} b) $\Longrightarrow$ a) (\footnote{uses W(1) only in weak form, without the saturation condition c). NB Only b) $\Longrightarrow$ c) uses W(2 bis), the implications a) $\Longrightarrow$ c) $\Longrightarrow$ b) are valid under the sole hypothesis W(1).}), for to see that $\pr_1 \in W$, applying the corollary above, we must prove that $\forall\ x \in X$
\[ (X \times \Dl^1)/x \fl X/x \tag*{$(*)$} \]
is in $W$. Now
\[ (X \times \Dl^1)/x \simeq X/x \times \Dl^1, \]
\[
\begin{tikzcd}
\Dl^1 \arrow[d] & X \times \Dl^1 \arrow[l] \arrow[d] & X/x \times \Dl^1 \arrow[l] \arrow[d] \\
e & X \arrow[l] & X/x \arrow[l]
\end{tikzcd}
\]
\texttt{[the vertical arrow on the right is added]} and the two categories $X/x$, $X/x \times \Dl^1$ are $W$-aspheric since they have a final object, so $(*)$ is in $W$ by weak saturation W(1).

a) $\Longrightarrow$ c). I recall that
\[ f : X \fl Y \]
\msp{page 2}
is a homotopism if there exists
\[ g : Y \fl X \]
such that
\[ gf \sim \id_X, \quad fg \sim \id_Y, \tag*{$(*)$} \]
where the relation $(\sim)$ denotes homotopy in $\uHom(X,X)$ and $\uHom(Y,Y)$ respectively (defined by paths going from one element to the other). Let us prove the

\textbf{Lemma.} \textit{Suppose that $W$ satisfies} W(1) \textit{and the} a) \textit{of the proposition, i.e. $X \times \Dl^1 \fl X$ in $W$ for every $X$ in} $\Cat$. \textit{Then if $X$, $Y$} \texttt{[are]} \textit{in} $\Cat$ \textit{and $f, g$ in} $\uHom(X,Y)$ \textit{are homotopic, we have}
\[ \hot_W(f) = \hot_W(g). \]
\textit{Moreover, $f \in W$ if and only if $g \in W$.}

It suffices to prove this when there exists a homomorphism $f \fl g$. This can be interpreted as the existence of an arrow
\[ X \times \Dl^1 \xrightarrow{F} Y \]
whose restrictions to $X \times \{0\}$ and $X \times \{1\}$ respectively are $f$, $g$. To prove that $\hot_W(f) = \hot_W(g)$, it suffices to prove that
\[ i_0, i_1 : X \fl X \times \Dl^1 \qquad x \longmapsto (x,0),\ x \longmapsto (x,1) \]
\printedpage{3}
% head: A. Grothendieck : Dérivateurs, Ch. XII, Édition des Universités de Montpellier II et Paris VII
% foot: Version 7 décembre 2002
are equal in $\Hot_W$, since
\[ f = F \circ i_0, \quad g = F \circ i_1. \]
Now this follows from the fact that they are sections of $X \times \Dl^1$ over $X$, and that in $\Hot_W$ the projection $X \times \Dl^1 \xrightarrow{p} X$ is \texttt{[an]} isomorphism (hence $i_0 = i_1 = p^{-1}$ in $\Hot_W$). Moreover, the diagram
\[
\begin{tikzcd}
X \arrow[r, "i_0"] \arrow[dr, "f"'] & X \times \Dl^1 \arrow[d, "F"] \\
& Y,
\end{tikzcd}
\]
where $i_0 \in W$, shows that $f \in W$ if and only if $F \in W$, and similarly $g \in W$ if and only if $F \in W$, hence $f \in W \Longleftrightarrow g \in W$.

Let us return to $(*)$, we see that in $\Hot_W$ $g \circ f = \id_X$, $f \circ g = \id_Y$, so $f$ and $g$ are $\Hot_W$-isomorphisms inverse to one another. On the other hand, $gf \in W$, $fg \in W$, since $\id_X, \id_Y \in W$. So we win, if we have formulated W(1) strongly enough. We have thus proved a) $\Longrightarrow$ c).

c) $\Longrightarrow$ b) is trivial, since if $X$ has a final object, $X \fl e$ is a homotopism.

\msp{page 3}

We must of course in what follows suppose that $W$ satisfies the equivalent conditions of the proposition. We shall take e.g. the form

\begin{itemize}[leftmargin=*]
\item[\textbf{W(3)}] (Final element axiom.) If $X$ has a final object, then $X$ is $W$-aspheric.
\end{itemize}

\textbf{Corollary 1.} \textit{For $f : X \fl Y$ to be $W$-aspheric, it is necessary and sufficient that the $X/y$ be $W$-aspheric} ($y \in \Ob Y$).

For in order that
\begin{tikzcd}[baseline=(\tikzcdmatrixname-1-1.base), column sep=small, row sep=small]
X/y \arrow[rr] \arrow[dr] & & Y/y \arrow[dl] \\
& e &
\end{tikzcd}
be in $W$, since $Y/y$ is $W$-aspheric (final object), it is necessary and sufficient that $X/y$ be so as well.

\textbf{Corollary 2.} \textit{If $X$ in} $\Cat$ \textit{has an initial object, it is $W$-aspheric.} (Indeed, $X \fl e$ is a homotopism, and we apply prop. 1.) (\footnote{This is thus the statement dual to the axiom W(3) of the final object.})

I shall put down for the record

\textbf{Corollary 3.} \textit{Every homotopism in} $\Cat$ \textit{is in $W$. If $f, g : X \rightrightarrows Y$ in} $\Cat$ \textit{are homotopic, then they are equal in} $\Hot_W$\textit{, and $f$ is in $W$ if and only if $g$ is.}

\textbf{Proposition 2.} Let $f : X \fl Y$ be in $\Cat$. If $f \in W$, $f^\circ \in W^\circ$ \texttt{[rather $f^\circ \in W$].}

\end{document}
